\addcontentsline{toc}{section}{CAPÍTULO V}

\chapter{AUTODIAGNÓSTICO INICIAL}

\section{Hipótesis de partida}

La hipótesis central que el MVP busca validar es la siguiente: si NexusCare se implementa como sistema de soporte a la auditoría de liquidaciones médicas, entonces entre el 50 \% y el 60 \% de las liquidaciones hospitalarias calificarán para el carril de \textit{fast-track}, reduciendo el tiempo total de auditoría por lote en al menos un 40 \%, porque el sistema clasifica las liquidaciones de bajo riesgo con justificaciones explicables que permiten al auditor validarlas en aproximadamente 1 minuto, y prioriza la cola de trabajo restante por nivel de riesgo clínico-económico.

Esta hipótesis se sustenta en tres hipótesis técnicas validadas durante Capstone 1:

\noindent 1. Los LLM pueden identificar inconsistencias clínicas en las liquidaciones hospitalarias. Validado parcialmente: MedGemma-4B generó análisis clínicos estructurados para el 93 \% de las liquidaciones (7 casos fallaron por truncamiento del contexto de 8,192 \textit{tokens}). La calidad clínica de los análisis exitosos no fue evaluada con métricas de precisión, dado que no existía un \textit{ground truth} de hallazgos esperados en Capstone 1. La validación completa de esta hipótesis se realizó en Capstone 2 mediante el benchmark de 5 modelos y la ejecución del pipeline integrado.

\noindent 2. El tiempo de procesamiento puede reducirse de 45 minutos a 5 minutos o menos por liquidación. Validado: el modelo genera análisis clínicos estructurados en menos de 30 segundos por caso, al que se suma el tiempo de revisión del auditor.

\noindent 3. Los puntajes de riesgo (0–100) con justificaciones claras son técnicamente viables. Validado: el modelo genera puntajes numéricos, acompañados de explicaciones clínicas textuales, que el auditor puede evaluar.

Sin embargo, Capstone 1 dejó abiertas preguntas sobre la calibración del modelo (tendencia a sobrevalorar el riesgo), la integración con las reglas contractuales específicas del plan y la protección de los datos personales en el flujo de procesamiento. Estas tres brechas definieron la

agenda técnica de Capstone 2: un \textit{benchmark} multimodelo para encontrar la calibración óptima, un motor de reglas contractuales que complementara al LLM y un módulo de anonimización que impidiera que datos personales llegaran al modelo de lenguaje.

\section{Supuestos no validados}

Al inicio de Capstone 2, tres supuestos permanecían sin validar Tabla~\ref{tab:supuestos-no-validados}
\begin{table}[H]
    \centering
    \caption{Supuestos no validados al inicio de Capstone 2}
    \label{tab:supuestos-no-validados}

    \begin{tabular}{
        p{0.37\textwidth}
        p{0.14\textwidth}
        p{0.41\textwidth}
    }
        \toprule
        \textbf{Supuesto}
        & \textbf{Criticidad}
        & \textbf{Consecuencia si es falso} \\
        \midrule

        Integración fluida con ERPs/TEDEF reales
        & Alta
        & Cuellos de botella técnicos impiden el procesamiento en tiempo real. \\

        Adopción y confianza del auditor médico
        & Alta
        & El auditor ignora las recomendaciones del sistema, anulando el beneficio operativo. \\

        Acceso a datos históricos para calibración
        & Media
        & Imposibilidad de ajustar los parámetros de riesgo del modelo. \\

        \bottomrule
    \end{tabular}
\end{table}

El primer supuesto se abordó parcialmente mediante la normalización de los datos crudos a un formato JSON estandarizado, simulando la capa de ingesta que existiría en producción. El segundo motivó el énfasis en la explicabilidad y la transparencia de las justificaciones, reconociendo que la tasa de \textit{fast-track} solo genera valor si el auditor confía en las recomendaciones que recibe. El tercero se mitigó mediante la construcción de un \textit{Golden Set}: un conjunto de 17 liquidaciones hospitalarias curadas y anotadas manualmente que sirve como referencia para evaluar la calidad de las recomendaciones del sistema.

\section{Matriz de riesgos} 
\begin{table}[H]
    \centering
    \caption{Matriz inicial de riesgos del proyecto}
    \label{tab:matriz-riesgos-inicial}

    \footnotesize
    \setlength{\tabcolsep}{4pt}
    \renewcommand{\arraystretch}{1.15}

    \begin{tabular}{
        p{0.21\textwidth}
        p{0.12\textwidth}
        p{0.10\textwidth}
        p{0.31\textwidth}
        p{0.18\textwidth}
    }
        \toprule
        \textbf{Riesgo}
        & \textbf{Probabilidad}
        & \textbf{Impacto}
        & \textbf{Estrategia de mitigación}
        & \textbf{Indicador de alerta} \\
        \midrule

        Baja adopción por desconfianza del auditor
        & Media
        & Alto
        & Exposición de la justificación clínica completa junto al
          \textit{score}; medición de la tasa de intervención.
        & Tasa de intervención $> 80$\,\%. \\

        Alucinaciones del modelo
        & Media
        & Alto
        & \textit{Human-in-the-loop} obligatorio; validación contra
          un \textit{Golden Set} de 17 casos curados.
        & Discordancia $> 30$\,\% con el \textit{Golden Set}. \\

        Calidad de los datos de integración
        & Alta
        & Medio
        & Validación estructural con Pydantic v2; monitoreo de la
          tasa de rechazo técnico.
        & Rechazo técnico $> 15$\,\%. \\

        Incumplimiento regulatorio (Ley N.\textsuperscript{o} 31814)
        & Baja
        & Alto
        & Supervisión humana obligatoria en todas las decisiones;
          trazabilidad completa de las justificaciones y
          anonimización de datos personales antes del procesamiento
          por el LLM.
        & Auditoría sin registro de validación humana. \\

        \bottomrule
    \end{tabular}
\end{table}


\newpage
